\documentclass[10pt,a4paper]{amsart}
\usepackage{amsmath,amssymb,amsthm,graphicx}
% Retained minimal preamble for the Stacks Project chapter more-algebra.tex (More on Algebra).
% The chapter file pinned-original-source.tex carries only \input{preamble} on its first line, so no
% preamble text lives inside the chapter itself.  The Stacks Project uses one single shared preamble for
% every chapter; the snapshot stages it identically as sources/stacks-chow/preamble.tex and
% sources/stacks-etale-cohomology/preamble.tex.  The block below is copied mechanically, line for line,
% from sources/stacks-chow/preamble.tex lines 14-18 (the Xy-pic setup) and 154-217 (the native theorem
% declarations together with the mathematical macros).  Nothing was added, removed or redefined; the
% wrapper class and the amsmath/amssymb/amsthm/graphicx base packages are supplied by the assembler.
% Chapter: more-algebra.tex -- More on Algebra.
% For commutative diagrams we use Xy-pic
\usepackage[all]{xy}

% We use 2cell for 2-commutative diagrams.
\xyoption{2cell}

% Theorem environments.
%
\theoremstyle{plain}
\newtheorem{theorem}[subsection]{Theorem}
\newtheorem{proposition}[subsection]{Proposition}
\newtheorem{lemma}[subsection]{Lemma}

\theoremstyle{definition}
\newtheorem{definition}[subsection]{Definition}
\newtheorem{example}[subsection]{Example}
\newtheorem{exercise}[subsection]{Exercise}
\newtheorem{situation}[subsection]{Situation}

\theoremstyle{remark}
\newtheorem{remark}[subsection]{Remark}
\newtheorem{remarks}[subsection]{Remarks}

\numberwithin{equation}{subsection}

% Macros
%
\def\lim{\mathop{\mathrm{lim}}\nolimits}
\def\colim{\mathop{\mathrm{colim}}\nolimits}
\def\Spec{\mathop{\mathrm{Spec}}}
\def\Hom{\mathop{\mathrm{Hom}}\nolimits}
\def\Ext{\mathop{\mathrm{Ext}}\nolimits}
\def\SheafHom{\mathop{\mathcal{H}\!\mathit{om}}\nolimits}
\def\SheafExt{\mathop{\mathcal{E}\!\mathit{xt}}\nolimits}
\def\Sch{\mathit{Sch}}
\def\Mor{\mathop{\mathrm{Mor}}\nolimits}
\def\Ob{\mathop{\mathrm{Ob}}\nolimits}
\def\Sh{\mathop{\mathit{Sh}}\nolimits}
\def\NL{\mathop{N\!L}\nolimits}
\def\CH{\mathop{\mathrm{CH}}\nolimits}
\def\proetale{{pro\text{-}\acute{e}tale}}
\def\etale{{\acute{e}tale}}
\def\QCoh{\mathit{QCoh}}
\def\Ker{\mathop{\mathrm{Ker}}}
\def\Im{\mathop{\mathrm{Im}}}
\def\Coker{\mathop{\mathrm{Coker}}}
\def\Coim{\mathop{\mathrm{Coim}}}

% Boxtimes
%
\DeclareMathSymbol{\boxtimes}{\mathbin}{AMSa}{"02}

%
% Macros for moduli stacks/spaces
%
\def\QCohstack{\mathcal{QC}\!\mathit{oh}}
\def\Cohstack{\mathcal{C}\!\mathit{oh}}
\def\Spacesstack{\mathcal{S}\!\mathit{paces}}
\def\Quotfunctor{\mathrm{Quot}}
\def\Hilbfunctor{\mathrm{Hilb}}
\def\Curvesstack{\mathcal{C}\!\mathit{urves}}
\def\Polarizedstack{\mathcal{P}\!\mathit{olarized}}
\def\Complexesstack{\mathcal{C}\!\mathit{omplexes}}
% \Pic is the operator that assigns to X its picard group, usage \Pic(X)
% \Picardstack_{X/B} denotes the Picard stack of X over B
% \Picardfunctor_{X/B} denotes the Picard functor of X over B
\def\Pic{\mathop{\mathrm{Pic}}\nolimits}
\def\Picardstack{\mathcal{P}\!\mathit{ic}}
\def\Picardfunctor{\mathrm{Pic}}
\def\Deformationcategory{\mathcal{D}\!\mathit{ef}}

% --- appended 2026-10-12 (Chief): the project own definitions of the reference/slogan/history
% display environments, copied mechanically from the shared Stacks preamble lines 8-12
% (\usepackage{verbatim} plus the three \newenvironment lines, which the chapter preamble keeps
% verbatim from the project preamble). Nothing added, removed or redefined: these environments
% carry no mathematical content and the project itself typesets them as comments. Without these
% three lines a counted statement whose own text uses them fails with Environment slogan undefined.
\usepackage{verbatim}
\newenvironment{reference}{\comment}{\endcomment}
\newenvironment{slogan}{\comment}{\endcomment}
\newenvironment{history}{\comment}{\endcomment}
\usepackage{cleveref}
\crefname{originallemma}{original lemma}{original lemmas}
\Crefname{originallemma}{Original lemma}{Original lemmas}
\makeatletter
\newlabel{algebra-lemma-valuation-ring-x-or-x-inverse}{{00IB}{source}{}{sourcepointer-1}{}}
\newlabel{algebra-lemma-valuation-ring-x-or-x-inverse@cref}{{[originallemma][0][]00IB}{[1][0][]source}{}{}{}}
\makeatother
\title{Rings with divisibility comparable for every pair of elements, characterized by principal finitely generated ideals with a unique maximal ideal}
\author{The Stacks Project Authors}
\date{stacks-project:more-algebra.tex:local-snapshot-20261012; GNU FDL 1.2 or later}
\begin{document}
\maketitle
\setcounter{subsection}{1}
\noindent\textit{Editorial note.} All delivered mathematical text is the author's own, copied from the pinned Stacks Project chapter file more-algebra.tex; no statement or proof step was rewritten, reordered or removed. The counted statement is the labeled lemma environment at source lines 37637-37649 (canonical tag 0ASN) and the counted proof is the proof environment at source lines 37651-37672, which immediately follows it and is its own proof.\par\medskip
\noindent\textit{Editorial note.} The counted statement keeps the author's own title Generalized valuation rings on the lemma environment and contains the author's reference environment recording the source of the result; the retained preamble defines that environment as comment-type, so its body and the citation inside it are discarded when the document is typeset, which is why the card needs no bibliography.\par\medskip
\noindent\textit{Editorial note.} The proof states that the equivalence of the first and third conditions is straightforward and does not spell it out; the two substantive implications, from the second condition to comparability and from comparability back to the second condition, are both carried out in full inside the counted span, so the card is not an instance of an omitted argument.\par\medskip
\noindent\textit{Editorial note.} The preamble span runs to line 90, the current end of the retained preamble, which is exactly what makes the reference environment a comment; the excerpt serves only as preamble, never as a span coordinate.\par\medskip
\noindent\textit{Editorial note.} The source URL is the source-level Stacks Project URL, as required when more than one card is produced from one source; the per-claim tag is carried by claim\_key and by this note.\par\medskip
\noindent\textit{Allowed standard prerequisites (editorial).}\begin{itemize}
\item In a local ring an element whose complement from a relation involving two generators lies in the maximal ideal forces one of the two coefficients to be a unit. used at proof lines 37656-37661. Produces the comparability of divisibility from the principality of the ideal generated by two elements.
\item Two distinct maximal ideals can be separated by elements lying in each but not the other, and such elements are mutually nondivisible. used at proof lines 37665-37669. Shows that comparability of divisibility forces the ring to be local.
\item If every two elements are comparable by divisibility then every finitely generated ideal is principal by induction, and a valuation ring is an example. used at proof lines 37669-37672. Completes the cycle of equivalences and handles the final statement.
\end{itemize}\medskip
\ifdefined\hypertarget\hypertarget{sourcepointer-1}{}\fi
\noindent\textit{Source reference.} Original lemma 00IB: sources/stacks-chow/tags.txt, source lines 673--673. The reference is cited at proof lines 37670-37672 for the final statement, that valuation rings satisfy the conditions of the lemma; it is a pointer to a standard characterization and the cycle of equivalences proved here is inside the counted span.\par\medskip
\begin{lemma}[Generalized valuation rings]
\label{lemma-generalized-valuation-ring}
\begin{reference}
\cite{Warfield-Decomposition}
\end{reference}
Let $R$ be a nonzero ring. The following are equivalent
\begin{enumerate}
\item For $a, b \in R$ either $a$ divides $b$ or $b$ divides $a$.
\item Every finitely generated ideal is principal and $R$ is local.
\item The set of ideals of $R$ is linearly ordered by inclusion.
\end{enumerate}
This holds in particular if $R$ is a valuation ring.
\end{lemma}

\begin{proof}
Assume (2) and let $a, b \in R$. Then $(a, b) = (c)$. If $c = 0$,
then $a = b = 0$ and $a$ divides $b$. Assume $c \not = 0$. Write
$c = ua + vb$ and $a = wc$ and $b = zc$. Then $c(1 - uw - vz) = 0$.
Since $R$ is local, this implies that $1 - uw - vz \in \mathfrak m$.
Hence either $w$ or $z$ is a unit, so either $a$ divides $b$ or
$b$ divides $a$. Thus (2) implies (1).

\medskip\noindent
Assume (1). If $R$ has two maximal ideals $\mathfrak m_i$
we can choose $a \in \mathfrak m_1$ with $a \not \in \mathfrak m_2$
and $b \in \mathfrak m_2$ with $b \not \in \mathfrak m_1$.
Then $a$ does not divide $b$ and $b$ does not divide $a$.
Hence $R$ has a unique maximal ideal and is local.
It follows easily from condition (1) and induction that every
finitely generated ideal is principal. Thus (1) implies (2).

\medskip\noindent
It is straightforward to prove that (1) and (3) are equivalent.
The final statement is Algebra, Lemma
\ref{algebra-lemma-valuation-ring-x-or-x-inverse}.
\end{proof}

\end{document}
